\section*{CHAPTER 4. CONCLUSION}
\addcontentsline{toc}{section}{\numberline{}CHAPTER 4. CONCLUSION}

\setcounter{section}{4}
\setcounter{subsection}{0}
\setcounter{figure}{0}
\setcounter{table}{0}
\setcounter{equation}{0}

%============================================================%
\subsection{Summary of the Study}

This report presented a comparative study between Orthogonal Frequency
Division Multiplexing (OFDM) and Single-Carrier Frequency Division
Multiple Access (SC-FDMA), two important transmission techniques used in
modern wireless communication systems.

Theoretical analysis showed that OFDM offers high spectral efficiency,
simple frequency-domain equalization, and strong robustness against
multipath fading. These advantages explain why OFDM has been widely
adopted in broadband wireless standards.

However, OFDM also suffers from high Peak-to-Average Power Ratio (PAPR),
which reduces power amplifier efficiency and increases the risk of
nonlinear distortion.

To address this limitation, SC-FDMA introduces a DFT precoding stage
before OFDM modulation. This preserves many OFDM advantages while
producing a lower-PAPR waveform with single-carrier characteristics.

%============================================================%
\subsection{Main Findings}

Based on theoretical discussion and simulation framework, the main
findings can be summarized as follows:

\begin{itemize}
    \item OFDM and SC-FDMA can achieve similar Bit Error Rate (BER)
    performance under identical modulation and channel conditions.

    \item SC-FDMA provides significantly lower PAPR than OFDM.

    \item Lower PAPR improves transmitter power efficiency and is highly
    beneficial for battery-powered mobile devices.

    \item This is the main reason SC-FDMA was selected for LTE uplink,
    while OFDM remained suitable for downlink transmission.
\end{itemize}

%============================================================%
\subsection{Practical Implications}

The comparison demonstrates that communication system design always
involves tradeoffs. OFDM emphasizes implementation simplicity and
multicarrier flexibility, whereas SC-FDMA improves power efficiency at
the transmitter.

Therefore, the preferred scheme depends on application requirements:

\begin{itemize}
    \item Base stations with less strict power constraints may favor OFDM.
    \item Mobile user equipment benefits from SC-FDMA.
    \item Future systems may combine advantages of both techniques.
\end{itemize}

%============================================================%
\subsection{Limitations of This Report}

Several simplifications were adopted in this study:

\begin{itemize}
    \item Only basic AWGN channel assumptions were considered.
    \item Channel coding effects were not included.
    \item Synchronization errors and hardware nonlinearities were not modeled.
    \item Simulation parameters were simplified for educational purposes.
\end{itemize}

More advanced models may produce deeper practical insights.

%============================================================%
\subsection{Future Work}

Possible future extensions include:

\begin{itemize}
    \item Comparison under Rayleigh and frequency-selective fading channels.
    \item Inclusion of channel coding schemes.
    \item Analysis of computational complexity.
    \item Investigation of PAPR reduction techniques for OFDM.
    \item Comparison with newer waveforms for 5G and beyond systems.
\end{itemize}

%============================================================%
\subsection{Final Remark}

OFDM and SC-FDMA are both elegant solutions to the challenge of reliable
high-speed wireless transmission. Their comparison illustrates a central
engineering principle: every improvement solves one problem while often
creating another. The best design is therefore not universal, but
context-dependent.